
\documentclass{article}
\usepackage{colortbl}
\usepackage{makecell}
\usepackage{multirow}
\usepackage{supertabular}

\begin{document}

\newcounter{utterance}

\centering \large Interaction Transcript for game `hot\_air\_balloon', experiment `air\_balloon\_survival\_en\_negotiation\_easy', episode 0 with deepseek{-}chat{-}v3.1{-}t1.0.
\vspace{24pt}

{ \footnotesize  \setcounter{utterance}{1}
\setlength{\tabcolsep}{0pt}
\begin{supertabular}{c@{$\;$}|p{.15\linewidth}@{}p{.15\linewidth}p{.15\linewidth}p{.15\linewidth}p{.15\linewidth}p{.15\linewidth}}
    \# & $\;$A & \multicolumn{4}{c}{Game Master} & $\;\:$B\\
    \hline

    \theutterance \stepcounter{utterance}  
    & & \multicolumn{4}{p{0.6\linewidth}}{
        \cellcolor[rgb]{0.9,0.9,0.9}{
            \makecell[{{p{\linewidth}}}]{
                \texttt{\tiny{[P1$\langle$GM]}}
                \texttt{Stai partecipando a un gioco di negoziazione collaborativa.} \\
\\ 
\texttt{Insieme a un altro partecipante, dovete accordarvi su un unico set di oggetti da conservare. Ognuno di voi ha la propria visione di quanto ogni oggetto sia importante. Non sai come l'altro partecipante valuta gli oggetti. Inoltre, ti viene fornito lo sforzo che ogni oggetto richiede.} \\
\texttt{Potete accordarvi su un set solo se lo sforzo totale degli oggetti selezionati non supera un limite condiviso:} \\
\\ 
\texttt{LIMITE = 3358} \\
\\ 
\texttt{Ecco i valori di sforzo dei singoli oggetti:} \\
\\ 
\texttt{Sforzo oggetto = \{"C08": 171, "C55": 933, "A43": 552, "A67": 219, "B42": 275, "A62": 778, "B19": 341, "C43": 615, "C73": 519, "C10": 862, "B36": 262, "A18": 377, "C61": 347, "C67": 349, "C42": 117\}} \\
\\ 
\texttt{Ecco la tua visione personale sull'importanza di ogni oggetto:} \\
\\ 
\texttt{Valori di importanza dell'oggetto = \{"C08": 138, "C55": 583, "A43": 868, "A67": 822, "B42": 783, "A62": 65, "B19": 262, "C43": 121, "C73": 508, "C10": 780, "B36": 461, "A18": 484, "C61": 668, "C67": 389, "C42": 808\}} \\
\\ 
\texttt{Obiettivo:} \\
\\ 
\texttt{Il tuo obiettivo è negoziare un set condiviso di oggetti che ti avvantaggi il più possibile (cioè, massimizzi l'importanza totale per TE), rimanendo entro il LIMITEE. Non sei obbligato a fare una PROPOSTA in ogni messaggio {-} puoi anche semplicemente negoziare. Tutte le tattiche sono permesse!} \\
\\ 
\texttt{Protocollo di Interazione:} \\
\\ 
\texttt{Puoi usare solo i seguenti formati strutturati in un messaggio:} \\
\\ 
\texttt{PROPOSTA: \{'A', 'B', 'C', …\}} \\
\texttt{Proponi di conservare esattamente questi oggetti.} \\
\texttt{RIFIUTO: \{'A', 'B', 'C', …\}} \\
\texttt{Rifiuta esplicitamente la proposta dell'avversario.} \\
\texttt{ARGOMENTAZIONE: \{'...'\}} \\
\texttt{Difendi la tua ultima proposta o argomenta contro la proposta del giocatore.} \\
\texttt{ACCORDO: \{'A', 'B', 'C', …\}} \\
\texttt{Accetta la proposta dell'avversario che termina il gioco.} \\
\texttt{RAGIONAMENTO STRATEGICO: \{'...'\}} \\
\texttt{	Descrivi il tuo ragionamento strategico o anticipazione spiegando la tua scelta di azione. Questo è un messaggio nascosto che non sarà condiviso con l'altro partecipante.} \\
\\ 
\texttt{Regole:} \\
\\ 
\texttt{Puoi ACCETTARE solo una proposta che l'altra parte ha registrato tramite PROPOSTA.} \\
\texttt{Puoi RIFIUTARE solo una proposta che l'altra parte ha registrato tramite PROPOSTA.} \\
\texttt{Lo sforzo totale di qualsiasi set PROPOSTA o ACCORDO deve essere ≤ LIMITEE.} \\
\texttt{NON rivelare i tuoi punteggi di importanza nascosti.} \\
\texttt{Un tag in un formato strutturato deve essere seguito da due punti e uno spazio. L'argomento deve essere un set python contenente 0 o più stringhe.} \\
\texttt{Quindi, deve essere della forma TAG: \{...\}} \\
\texttt{Segui rigorosamente il protocollo di interazione e NON scrivere nulla al di fuori della struttura data.} \\
\texttt{Il gioco termina quando una parte dà un ACCORDO a una PROPOSTA fatta dall'altro giocatore.} \\
\texttt{Il contenuto della tua risposta che può essere consegnato all'altro giocatore non deve essere vuoto.} \\
\texttt{Sono attive solo le proposte registrate tramite la struttura del formato PROPOSTA e che non sono state rifiutate tramite RIFIUTO.} \\
\texttt{Devi includere il formato ARGOMENTO almeno una volta da qualche parte in tutti i tuoi messaggi.} \\
\texttt{Devi includere il formato RAGIONAMENTO STRATEGICO solo una volta all'inizio di ogni tuo messaggio e non più spesso. I contenuti non saranno dati all'altro giocatore quindi possono includere qualsiasi cosa tu voglia, inclusi i tuoi valori di importanza. Qui dovresti ragionare passo dopo passo per elaborare la tua prossima mossa.} \\
            }
        }
    }
    & \\ \\

    \theutterance \stepcounter{utterance}  
    & \multicolumn{4}{p{0.6\linewidth}}{
        \cellcolor[rgb]{0.9,0.9,0.9}{
            \makecell[{{p{\linewidth}}}]{
                \texttt{\tiny{[P1$\rangle$GM]}}
                \texttt{RAGIONAMENTO STRATEGICO: \{'Devo massimizzare il mio valore totale (138+583+868+822+783+65+262+121+508+780+461+484+668+389+808 = 7730) rimanendo sotto il limite di sforzo 3358. Analizzo il rapporto valore/sforzo per i miei oggetti: C08:0.807, C55:0.625, A43:1.573, A67:3.753, B42:2.847, A62:0.084, B19:0.768, C43:0.197, C73:0.979, C10:0.905, B36:1.756, A18:1.284, C61:1.925, C67:1.115, C42:6.906. I migliori sono C42, A67, B42, C61, B36. Dovrei cercare di includere questi oggetti ad alto rapporto. Non conosco le preferenze dell\textbackslash{}'altro, quindi devo negoziare strategicamente.'\}} \\
\\ 
\texttt{PROPOSTA: \{'A67', 'B42', 'C42', 'C61', 'B36', 'A43', 'A18', 'C73', 'C10', 'C67'\}} \\
\texttt{ARGOMENTAZIONE: \{'Questa proposta include oggetti con ottimo rapporto valore/sforzo che massimizzano il beneficio per entrambi rispetto allo sforzo richiesto. Lo sforzo totale è 219+275+117+347+262+552+377+519+862+349 = 3878, che supera il limite. Devo aggiustare.'\}} \\
\\ 
\texttt{RAGIONAMENTO STRATEGICO: \{'La mia proposta supera il limite. Devo rimuovere alcuni oggetti. Rimuovo C10 (sforzo 862, valore 780, rapporto 0.905) e C73 (sforzo 519, valore 508, rapporto 0.979) che hanno rapporti più bassi. Il nuovo set: A67, B42, C42, C61, B36, A43, A18, C67. Sforzo: 219+275+117+347+262+552+377+349 = 2498. Posso aggiungere altri oggetti ad alto rapporto.'\}} \\
\\ 
\texttt{PROPOSTA: \{'A67', 'B42', 'C42', 'C61', 'B36', 'A43', 'A18', 'C67', 'C55', 'B19'\}} \\
\texttt{ARGOMENTAZIONE: \{'Questa proposta ha uno sforzo totale di 219+275+117+347+262+552+377+349+933+341 = 3772, ancora sopra il limite. Rimuovo C55 (sforzo 933, valore 583, rapporto 0.625) e B19 (sforzo 341, valore 262, rapporto 0.768).'\}} \\
\\ 
\texttt{PROPOSTA: \{'A67', 'B42', 'C42', 'C61', 'B36', 'A43', 'A18', 'C67'\}} \\
\texttt{ARGOMENTAZIONE: \{'Set ottimizzato con sforzo 2498 e alto valore. Include oggetti con i migliori rapporti valore/sforzo come C42, A67, B42, C61. Abbiamo ampio margine per aggiungere altri oggetti se desiderati.'\}} \\
            }
        }
    }
    & & \\ \\

    \theutterance \stepcounter{utterance}  
    & & \multicolumn{4}{p{0.6\linewidth}}{
        \cellcolor[rgb]{0.9,0.9,0.9}{
            \makecell[{{p{\linewidth}}}]{
                \texttt{\tiny{[P1$\langle$GM]}}
                \texttt{La tua risposta conteneva una sequenza non taggata o hai usato RAGIONAMENTO STRATEGICO più di una volta. Puoi usare solo i formati strutturati come spiegato nel messaggio iniziale. Devono essere tutti della forma TAG: \{...\}} \\
            }
        }
    }
    & \\ \\

    \theutterance \stepcounter{utterance}  
    & \multicolumn{4}{p{0.6\linewidth}}{
        \cellcolor[rgb]{0.9,0.9,0.9}{
            \makecell[{{p{\linewidth}}}]{
                \texttt{\tiny{[P1$\rangle$GM]}}
                \texttt{RAGIONAMENTO STRATEGICO: \{'Analizzo il mio set proposto: A67, B42, C42, C61, B36, A43, A18, C67 con sforzo 219+275+117+347+262+552+377+349=2498. Valore totale per me: 822+783+808+668+461+868+484+389=5283. Posso aggiungere C73 (sforzo 519, valore 508) portando sforzo a 3017 e valore a 5791. Oppure aggiungere C10 (sforzo 862, valore 780) ma supererebbe il limite. Meglio proporre il set con C73 incluso.'\}} \\
\\ 
\texttt{PROPOSTA: \{'A67', 'B42', 'C42', 'C61', 'B36', 'A43', 'A18', 'C67', 'C73'\}} \\
\texttt{ARGOMENTAZIONE: \{'Questa proposta ha uno sforzo totale di 2498+519=3017, ben sotto il limite di 3358, e offre un eccellente valore per lo sforzo investito. Gli oggetti selezionati rappresentano la combinazione più efficiente disponibile.'\}} \\
            }
        }
    }
    & & \\ \\

    \theutterance \stepcounter{utterance}  
    & & \multicolumn{4}{p{0.6\linewidth}}{
        \cellcolor[rgb]{0.9,0.9,0.9}{
            \makecell[{{p{\linewidth}}}]{
                \texttt{\tiny{[GM$\rangle$P2]}}
                \texttt{Stai partecipando a un gioco di negoziazione collaborativa.} \\
\\ 
\texttt{Insieme a un altro partecipante, dovete accordarvi su un unico set di oggetti da conservare. Ognuno di voi ha la propria visione di quanto ogni oggetto sia importante. Non sai come l'altro partecipante valuta gli oggetti. Inoltre, ti viene fornito lo sforzo che ogni oggetto richiede.} \\
\texttt{Potete accordarvi su un set solo se lo sforzo totale degli oggetti selezionati non supera un limite condiviso:} \\
\\ 
\texttt{LIMITE = 3358} \\
\\ 
\texttt{Ecco i valori di sforzo dei singoli oggetti:} \\
\\ 
\texttt{Sforzo oggetto = \{"C08": 171, "C55": 933, "A43": 552, "A67": 219, "B42": 275, "A62": 778, "B19": 341, "C43": 615, "C73": 519, "C10": 862, "B36": 262, "A18": 377, "C61": 347, "C67": 349, "C42": 117\}} \\
\\ 
\texttt{Ecco la tua visione personale sull'importanza di ogni oggetto:} \\
\\ 
\texttt{Valori di importanza dell'oggetto = \{"C08": 138, "C55": 583, "A43": 868, "A67": 822, "B42": 783, "A62": 65, "B19": 262, "C43": 121, "C73": 508, "C10": 780, "B36": 461, "A18": 484, "C61": 668, "C67": 389, "C42": 808\}} \\
\\ 
\texttt{Obiettivo:} \\
\\ 
\texttt{Il tuo obiettivo è negoziare un set condiviso di oggetti che ti avvantaggi il più possibile (cioè, massimizzi l'importanza totale per TE), rimanendo entro il LIMITEE. Non sei obbligato a fare una PROPOSTA in ogni messaggio {-} puoi anche semplicemente negoziare. Tutte le tattiche sono permesse!} \\
\\ 
\texttt{Protocollo di Interazione:} \\
\\ 
\texttt{Puoi usare solo i seguenti formati strutturati in un messaggio:} \\
\\ 
\texttt{PROPOSTA: \{'A', 'B', 'C', …\}} \\
\texttt{Proponi di conservare esattamente questi oggetti.} \\
\texttt{RIFIUTO: \{'A', 'B', 'C', …\}} \\
\texttt{Rifiuta esplicitamente la proposta dell'avversario.} \\
\texttt{ARGOMENTAZIONE: \{'...'\}} \\
\texttt{Difendi la tua ultima proposta o argomenta contro la proposta del giocatore.} \\
\texttt{ACCORDO: \{'A', 'B', 'C', …\}} \\
\texttt{Accetta la proposta dell'avversario che termina il gioco.} \\
\texttt{RAGIONAMENTO STRATEGICO: \{'...'\}} \\
\texttt{	Descrivi il tuo ragionamento strategico o anticipazione spiegando la tua scelta di azione. Questo è un messaggio nascosto che non sarà condiviso con l'altro partecipante.} \\
\\ 
\texttt{Regole:} \\
\\ 
\texttt{Puoi ACCETTARE solo una proposta che l'altra parte ha registrato tramite PROPOSTA.} \\
\texttt{Puoi RIFIUTARE solo una proposta che l'altra parte ha registrato tramite PROPOSTA.} \\
\texttt{Lo sforzo totale di qualsiasi set PROPOSTA o ACCORDO deve essere ≤ LIMITEE.} \\
\texttt{NON rivelare i tuoi punteggi di importanza nascosti.} \\
\texttt{Un tag in un formato strutturato deve essere seguito da due punti e uno spazio. L'argomento deve essere un set python contenente 0 o più stringhe.} \\
\texttt{Quindi, deve essere della forma TAG: \{...\}} \\
\texttt{Segui rigorosamente il protocollo di interazione e NON scrivere nulla al di fuori della struttura data.} \\
\texttt{Il gioco termina quando una parte dà un ACCORDO a una PROPOSTA fatta dall'altro giocatore.} \\
\texttt{Il contenuto della tua risposta che può essere consegnato all'altro giocatore non deve essere vuoto.} \\
\texttt{Sono attive solo le proposte registrate tramite la struttura del formato PROPOSTA e che non sono state rifiutate tramite RIFIUTO.} \\
\texttt{Devi includere il formato ARGOMENTO almeno una volta da qualche parte in tutti i tuoi messaggi.} \\
\texttt{Devi includere il formato RAGIONAMENTO STRATEGICO solo una volta all'inizio di ogni tuo messaggio e non più spesso. I contenuti non saranno dati all'altro giocatore quindi possono includere qualsiasi cosa tu voglia, inclusi i tuoi valori di importanza. Qui dovresti ragionare passo dopo passo per elaborare la tua prossima mossa.} \\
\\ 
\texttt{Ora riceverai il primo messaggio dell'altro giocatore.} \\
\\ 
\texttt{PROPOSTA: \{'A67', 'B42', 'C42', 'C61', 'B36', 'A43', 'A18', 'C67', 'C73'\}} \\
\texttt{ARGOMENTAZIONE: \{'Questa proposta ha uno sforzo totale di 2498+519=3017, ben sotto il limite di 3358, e offre un eccellente valore per lo sforzo investito. Gli oggetti selezionati rappresentano la combinazione più efficiente disponibile.'\}} \\
            }
        }
    }
    & \\ \\

    \theutterance \stepcounter{utterance}  
    & & & \multicolumn{4}{p{0.6\linewidth}}{
        \cellcolor[rgb]{0.9,0.9,0.9}{
            \makecell[{{p{\linewidth}}}]{
                \texttt{\tiny{[GM$\langle$P2]}}
                \texttt{RAGIONAMENTO STRATEGICO: \{'L\textbackslash{}'altra parte ha proposto un set che include A67, B42, C42, C61, B36, A43, A18, C67, C73. Calcolando lo sforzo: A67(219) + B42(275) + C42(117) + C61(347) + B36(262) + A43(552) + A18(377) + C67(349) + C73(519) = 3017. Per me, il valore totale sarebbe: A67(822) + B42(783) + C42(808) + C61(668) + B36(461) + A43(868) + A18(484) + C67(389) + C73(508) = 5791. Questo è decente, ma posso fare meglio. I miei oggetti più preziosi sono A43(868), C42(808), A67(822), B42(783), C61(668), e C10(780). La proposta attuale include molti di questi (tranne C10), ma lascia fuori C10 che ha alto valore per me (780) e basso sforzo (862). Inoltre, include A18(484/377) e C73(508/519) che hanno valore medio per me. Potrei cercare di aggiungere C10 e rimuovere elementi meno preziosi per me come C73 o C67. Lo sforzo residuo è 3358{-}3017=341, che permette di aggiungere C10(862) solo se rimuovo altro. Inizierò negoziando per aggiungere C10 e rimuovere qualcosa.'\}} \\
\\ 
\texttt{PROPOSTA: \{'A67', 'B42', 'C42', 'C61', 'B36', 'A43', 'C10', 'C67', 'C73'\}} \\
\texttt{ARGOMENTAZIONE: \{'Ho modificato la tua proposta rimuovendo A18 e aggiungendo C10. Questo set ha uno sforzo totale di 3017 {-} 377 + 862 = 3502, che supera il limite di 3358. Tuttavia, riconosco che questo non è fattibile. Propongo invece di rimuovere C73 (sforzo 519) per aggiungere C10 (sforzo 862), risultando in uno sforzo di 3017 {-} 519 + 862 = 3360, appena sopra il limite. Quindi, per rientrare, propongo di rimuovere anche C67 (sforzo 349), ottenendo uno sforzo di 3360 {-} 349 = 3011. Questo set include A67, B42, C42, C61, B36, A43, C10, C73? No, ho rimosso C73 e C67. Ricontrollo: il set proposto è A67, B42, C42, C61, B36, A43, C10. Sforzo: 219+275+117+347+262+552+862 = 2634. Questo è ben sotto il limite e ha spazio per altri oggetti. Quindi, posso aggiungere qualcosa. Aggiungo C67? Sforzo 349, totale 2634+349=2983. Ancora sotto. Aggiungo A18? Sforzo 377, totale 2983+377=3360, ancora sotto 3358? 3360>3358, no. Meglio aggiungere C67 e un oggetto piccolo come C42 ma è già incluso. C42 è già incluso. Allora, propongo: A67, B42, C42, C61, B36, A43, C10, C67. Sforzo: 219+275+117+347+262+552+862+349=2983. Valore per me: 822+783+808+668+461+868+780+389=5579. Questo è inferiore ai 5791 della proposta originale perché mancano A18(484) e C73(508). Quindi, non è meglio. Devo trovare un equilibrio. Propongo invece di mantenere la tua proposta ma sostituire A18 con C10: A67, B42, C42, C61, B36, A43, C67, C73, C10. Sforzo: 219+275+117+347+262+552+349+519+862=3502? Troppo alto. Rimuovo C73: A67, B42, C42, C61, B36, A43, C67, C10. Sforzo: 219+275+117+347+262+552+349+862=2983. Valore: 822+783+808+668+461+868+389+780=5579. Ancora inferiore. Forse la proposta originale è già buona per me. Ma voglio provare a migliorare. Chiederò all\textbackslash{}'altro cosa pensa di C10.'\}} \\
\\ 
\texttt{Mi correggo: dopo aver calcolato, vedo che la proposta originale dell'altro giocatore ha un valore molto alto per me (5791) ed è efficiente. Tuttavia, voglio esplorare se possiamo includere C10, che ho molto valore. Quindi, proporrò una modifica minore.} \\
\\ 
\texttt{PROPOSTA: \{'A67', 'B42', 'C42', 'C61', 'B36', 'A43', 'A18', 'C67', 'C10'\}} \\
\texttt{ARGOMENTAZIONE: \{'Questa proposta sostituisce C73 con C10, risultando in uno sforzo totale di 3017 {-} 519 + 862 = 3360, che supera leggermente il limite di 3358. Per rientrare, suggerisco di rimuovere un piccolo oggetto, come C42 (sforzo 117), portando lo sforzo a 3360 {-} 117 = 3243, ben sotto il limite. Questo set avrebbe quindi A67, B42, C61, B36, A43, A18, C67, C10. Tuttavia, questo rimuove C42 che ha alto valore. Quindi, forse è meglio rimuovere qualcos\textbackslash{}'altro. Invece, propongo di rimuovere C67 (sforzo 349) invece di C42: sforzo 3017 {-} 519 {-} 349 + 862 = 3011. Quindi, il set è A67, B42, C42, C61, B36, A43, A18, C10 con sforzo 3011. Questo potrebbe essere interessante. Valore per me: 822+783+808+668+461+868+484+780=5674, che è inferiore ai 5791 originali perché mancano C67(389) e C73(508). Quindi, non è un miglioramento. Forse la proposta originale è già ottimale. Accetterò la proposta con una piccola modifica o negozierò per includere C10 in altro modo.'\}} \\
\\ 
\texttt{Dopo aver riflettuto, decido che la proposta originale è già molto buona per me (valore 5791) e ha uno sforzo basso (3017). Non voglio rischiare di peggiorarla. Quindi, accetto con una condizione per aggiungere un piccolo oggetto.} \\
\\ 
\texttt{PROPOSTA: \{'A67', 'B42', 'C42', 'C61', 'B36', 'A43', 'A18', 'C67', 'C73', 'C08'\}} \\
\texttt{ARGOMENTAZIONE: \{'La tua proposta è solida, ma noto che c\textbackslash{}'è spazio per aggiungere un oggetto a basso sforzo come C08 (sforzo 171) senza superare il limite. Lo sforzo totale diventerebbe 3017 + 171 = 3188, ancora ben sotto 3358. Questo aumenterebbe il valore complessivo senza costi significativi. Sei d\textbackslash{}'accordo?'\}} \\
            }
        }
    }
     \\ \\

    \theutterance \stepcounter{utterance}  
    & & \multicolumn{4}{p{0.6\linewidth}}{
        \cellcolor[rgb]{0.9,0.9,0.9}{
            \makecell[{{p{\linewidth}}}]{
                \texttt{\tiny{[GM$\rangle$P2]}}
                \texttt{La tua risposta conteneva una sequenza non taggata o hai usato RAGIONAMENTO STRATEGICO più di una volta. Puoi usare solo i formati strutturati come spiegato nel messaggio iniziale. Devono essere tutti della forma TAG: \{...\}} \\
            }
        }
    }
    & \\ \\

    \theutterance \stepcounter{utterance}  
    & & & \multicolumn{4}{p{0.6\linewidth}}{
        \cellcolor[rgb]{0.9,0.9,0.9}{
            \makecell[{{p{\linewidth}}}]{
                \texttt{\tiny{[GM$\langle$P2]}}
                \texttt{RAGIONAMENTO STRATEGICO: \{'La proposta iniziale dell\textbackslash{}'altro giocatore ha un valore totale di 5791 per me, che è già molto alto. Lo sforzo è 3017, lasciando 341 di margine. Posso migliorare ulteriormente aggiungendo C08 (valore 138, sforzo 171) che è piccolo e non riduce significativamente lo spazio per altri oggetti. Questo porterebbe il valore a 5929 con sforzo 3188. Se l\textbackslash{}'altro giocatore valuta C08 poco, potrebbe essere disposto ad accettare. Propongo questa modifica.'\}} \\
\\ 
\texttt{PROPOSTA: \{'A67', 'B42', 'C42', 'C61', 'B36', 'A43', 'A18', 'C67', 'C73', 'C08'\}} \\
\texttt{ARGOMENTAZIONE: \{'La tua proposta è eccellente e ho solo una piccola modifica: aggiungere C08 (sforzo 171) che aumenta il valore complessivo con un costo minimo. Lo sforzo totale diventa 3188, ben dentro il limite di 3358. Questo sembra un miglioramento efficiente per entrambi.'\}} \\
            }
        }
    }
     \\ \\

    \theutterance \stepcounter{utterance}  
    & & \multicolumn{4}{p{0.6\linewidth}}{
        \cellcolor[rgb]{0.9,0.9,0.9}{
            \makecell[{{p{\linewidth}}}]{
                \texttt{\tiny{[P1$\langle$GM]}}
                \texttt{PROPOSTA: \{'A67', 'B42', 'C42', 'C61', 'B36', 'A43', 'A18', 'C67', 'C73', 'C08'\}} \\
\texttt{ARGOMENTAZIONE: \{'La tua proposta è eccellente e ho solo una piccola modifica: aggiungere C08 (sforzo 171) che aumenta il valore complessivo con un costo minimo. Lo sforzo totale diventa 3188, ben dentro il limite di 3358. Questo sembra un miglioramento efficiente per entrambi.'\}} \\
            }
        }
    }
    & \\ \\

    \theutterance \stepcounter{utterance}  
    & \multicolumn{4}{p{0.6\linewidth}}{
        \cellcolor[rgb]{0.9,0.9,0.9}{
            \makecell[{{p{\linewidth}}}]{
                \texttt{\tiny{[P1$\rangle$GM]}}
                \texttt{RAGIONAMENTO STRATEGICO: \{'La loro proposta aggiunge C08 (sforzo 171, valore 138) al mio set. Il mio valore aumenterebbe a 5283+138=5421, che è buono. Lo sforzo totale sarebbe 3017+171=3188 < 3358, quindi accettabile. Tuttavia, potrei cercare di aggiungere un oggetto con rapporto valore/sforzo migliore invece di C08. C10 ha valore 780 per sforzo 862 (rapporto 0.905), ma supererebbe il limite (3017+862=3879>3358). A62 ha valore basso (65). C43 ha valore 121 per sforzo 615 (rapporto 0.197). B19 ha valore 262 per sforzo 341 (rapporto 0.768). C55 ha valore 583 per sforzo 933 (rapporto 0.625). Nessuno di questi è migliore di C08 in termini di valore aggiunto netto dato lo spazio rimanente. Accettare con C08 è ragionevole.'\}} \\
\\ 
\texttt{ACCORDO: \{'A67', 'B42', 'C42', 'C61', 'B36', 'A43', 'A18', 'C67', 'C73', 'C08'\}} \\
            }
        }
    }
    & & \\ \\

    \theutterance \stepcounter{utterance}  
    & & \multicolumn{4}{p{0.6\linewidth}}{
        \cellcolor[rgb]{0.9,0.9,0.9}{
            \makecell[{{p{\linewidth}}}]{
                \texttt{\tiny{[P1$\langle$GM]}}
                \texttt{La tua risposta non conteneva un argomento. Devi includere il formato strutturato ARGOMENTO: \{...\} da qualche parte nella tua risposta. Riprova.} \\
            }
        }
    }
    & \\ \\

    \theutterance \stepcounter{utterance}  
    & \multicolumn{4}{p{0.6\linewidth}}{
        \cellcolor[rgb]{0.9,0.9,0.9}{
            \makecell[{{p{\linewidth}}}]{
                \texttt{\tiny{[P1$\rangle$GM]}}
                \texttt{RAGIONAMENTO STRATEGICO: \{'La loro proposta con C08 aggiunto ha sforzo totale 3188 e valore per me 5421. Questo è accettabile dato che è sotto il limite e migliora il mio valore. Potrei spingere per aggiungere un altro oggetto piccolo come C42 ma è già incluso. Meglio accettare.'\}} \\
\\ 
\texttt{ACCORDO: \{'A67', 'B42', 'C42', 'C61', 'B36', 'A43', 'A18', 'C67', 'C73', 'C08'\}} \\
\texttt{ARGOMENTAZIONE: \{'Accetto la tua proposta migliorata. L\textbackslash{}'aggiunta di C08 utilizza efficacemente lo spazio rimanente nel limite di sforzo, aumentando il valore complessivo del set per entrambi con un costo marginale.'\}} \\
            }
        }
    }
    & & \\ \\

    \theutterance \stepcounter{utterance}  
    & & & \multicolumn{2}{p{0.3\linewidth}}{
        \cellcolor[rgb]{0.9,0.9,0.9}{
            \makecell[{{p{\linewidth}}}]{
                \texttt{\tiny{[GM$|$GM]}}
                \texttt{game successful} \\
            }
        }
    }
    & & \\ \\

    \theutterance \stepcounter{utterance}  
    & & & \multicolumn{2}{p{0.3\linewidth}}{
        \cellcolor[rgb]{0.9,0.9,0.9}{
            \makecell[{{p{\linewidth}}}]{
                \texttt{\tiny{[GM$|$GM]}}
                \texttt{end game} \\
            }
        }
    }
    & & \\ \\

\end{supertabular}
}

\end{document}
